\section{Related Work}
\label{sec:related}

We organise prior work along the eight themes that emerged from a systematic
sweep of the literature, and position our integration against each. Seven
differentiators recur: (A) UK FRA / Building-Safety-Act integration; (B)
Bayesian latent-state inference of building condition; (C) calibrated
uncertainty with missing-evidence flags; (D) inspection as noisy observation;
(E) rare-event outcome treatment; (F) a mechanism-informed consequence
surrogate; and (G) decision-theoretic inspection prioritisation.

\paragraph{UK fire-safety landscape and FRA methodology.}
A recent systematic review of UK building fire-safety practice catalogues its
technological, mitigation, behavioural, and regulatory strands, and flags
persistent compliance, enforcement, and maintenance gaps \citep{dauda2025uk};
the Hackitt review diagnosed a ``broken'' regulatory system and seeded the
outcomes-based safety-case regime that followed \citep{hackitt2018,perry2025}.
Quantitative fire risk assessment itself has a long lineage --- event-tree and
deterministic consequence analysis in textbook form \citep{yung2008,rasbash2004},
statistics-driven residential risk indices \citep{xin2013}, single-factor
effectiveness estimates such as smoke-alarm impact \citep{gilbert2021}, and
UK-specific trend analyses of accidental dwelling fires \citep{taylor2026}.
This literature motivates our setting (A) but is either policy diagnosis or
static, deterministic assessment; what an ``evidence-based'' regime lacks, and
what we supply, is the probabilistic latent-state machinery (B, C).

\paragraph{Probabilistic and Bayesian-network fire-risk models.}
The closest prior line chains fire development, escape, and intervention for UK
dwellings in three-part Bayesian networks (BNs)
\citep{matellini2018threepart,matellini2013}, alongside BN fatality and spread
models \citep{hanea2009,cheng2009} and a credal-network probabilistic risk
assessment of fire occurrence applied to Grenfell Tower
\citep{estradalugo2019grenfell}. More recent variants extend the construction
to heritage buildings \citep{chen2023hist}, subway and tunnel scenario
evolution \citep{li2024subway,hidayat2025tunnel}, and high-rise electrical and
fuzzy composite scoring \citep{su2023electrical,wang2021fuzzy,tan2021tho}, all
on the BN foundations of \citet{pearl1988}. We adopt this tradition's
latent-state, posterior view of risk --- but throughout it, conditional
probability tables (CPTs) are fixed from expert judgement or aggregate
statistics, inputs are treated as certain, and outputs are points: no
calibrated uncertainty (C), no treatment of inspections as noisy observations
(D), no physics-based consequence layer (F), and no inspection-prioritisation
layer (G).

\paragraph{Uncertainty quantification and decision-theoretic inspection.}
Risk matrices --- the format operational FRA actually uses --- have poor
resolution and can rank risks worse than random \citep{cox2008,thomas2014},
while calibration and sharpness under proper scoring rules provide the
evaluation standard for probabilistic forecasts that we adopt
\citep{gneiting2007}. Separately, the structural-reliability
risk-based-inspection and value-of-information tradition treats inspections as
imperfect information that updates a reliability estimate, and optimises
inspection planning on that basis
\citep{straub2005,luque2019,zhang2021voi,vereecken2020,klerk2019}. Our
contribution here is a transplant: importing inspection-as-noisy-observation
(D) and expected value of inspection (G) into fire risk assessment, and
replacing categorical scores with posterior distributions (C) --- a combination
this tradition has never instantiated for fire.

\paragraph{Fire-risk datasets, geospatial mapping, and service targeting.}
England-wide small-area evidence links deprivation to fire-related dwelling
casualties \citep{li2022}; emergency-service demand analytics
\citep{londonblue2025} and evaluations of home-fire-safety-visit targeting
\citep{waring2024} supply the operational context, with hierarchical Bayesian
small-area rate estimation as the methodological template \citep{macnab2003}.
The open-datacube precedent is IberFire, a spatio-temporal wildfire-risk cube
assembled entirely from open sources \citep{erzibengoa2025iberfire}. These
works supply our covariates and the open-data blueprint (A), but remain
ecological and associational: none couples the data to a building-conditioned
probabilistic model with quantified uncertainty (B, C).

\paragraph{Machine learning in fire safety.}
Reviews of fire-safety machine learning map a fast-growing field and flag
precisely the weaknesses we target: underdeveloped uncertainty quantification
and calibration \citep{yildiz2025ml,tapehnaser2022}, with mechanistically grounded
learning advocated as the remedy \citep{naser2021}. Applications span
smart-building design \citep{zeng2024smart}, classifiers explained post hoc
with SHapley additive explanations (SHAP) \citep{lu2023stadium,dai2025interp},
and knowledge-graph and city-scale forecasting
\citep{xu2025knowgraph,lee2025urbanfire}. We use machine learning differently:
as a calibrated surrogate \emph{inside} a probabilistic model (C, F), rather
than as a black-box or post-hoc-explained point classifier.

\paragraph{Fast fire-simulation surrogates.}
A systematic review of physics-informed surrogates finds them rarely embedded
inside a probabilistic risk model \citep{yarmohammadian2025pism} --- the
integration gap this paper occupies. The component technology is mature:
emulators predict fire fields, flashover, and detection orders of magnitude
faster than computational fluid dynamics
\citep{lattimer2020,wang2021pflash,nguyen2022,zeng2022ifetool,zeng2024jcde};
physics-informed neural networks learn interpretable spread parameters
\citep{vogiatzoglou2025}; probabilistic diffusion surrogates output spread
ensembles \citep{yu2026}; and graph neural networks serve as surrogates over
building-block graphs \citep{wu2025gnn}. We adopt the surrogate-on-graph idea
for our archetypes (F) but couple it into the risk posterior with propagated
uncertainty (C).

\paragraph{Coupled fire--evacuation modelling.}
The deterministic available-versus-required safe-egress-time (ASET/RSET) margin
has been criticised as conceptually flawed \citep{babrauskas2010} and reframed
distributionally \citep{schroeder2020}; one-way fire-to-egress coupling and
agent-based evacuation models degrade movement under tenability loss
\citep{cheong2021,wang2025medbuild}, with occupant vulnerability shaping
evacuation outcomes \citep{tancogne2016}. We follow the distributional turn to
its conclusion: evacuation failure is a node inferred \emph{within} the risk
model with propagated uncertainty (C, E), not a deterministic single-scenario
margin.

\paragraph{Rare-event statistics.}
Extreme-value and peaks-over-threshold methods model heavy-tailed fire losses
\citep{ramachandran1974,mcneil1997,embrechts1997}, and survival models give an
interpretable time-to-event template \citep{kalair2021}. We apply this
machinery to rare-event fire outcomes coupled to a building-conditioned
posterior (E) --- a link absent from its actuarial and traffic source domains.

\paragraph{Positioning.}
Our closest competitors, in order: the UK dwelling-fire Bayesian-network line
\citep{matellini2018threepart,matellini2013}, which shares our setting and
graphical structure but fixes its CPTs, treats inputs as certain, and offers
neither calibrated intervals nor a latent-state posterior, inspection-noise
model, physics surrogate, or prioritisation layer; the structural
risk-based-inspection tradition \citep{straub2005,luque2019}, the closest
analogue to our contributions (D) and (G) but developed for structural
deterioration and never instantiated for fire; the physics-informed surrogate
line \citep{yarmohammadian2025pism,naser2021}, which defines the surrogate half
of (F) but almost never embeds it in a risk model with quantified uncertainty;
and the risk-matrix critique \citep{cox2008,thomas2014}, which justifies
abandoning categorical outputs (C) but supplies no constructive replacement.
Every competitor covers one or two of the seven differentiators; none couples
UK regulatory context, a Bayesian latent-state model of building condition,
calibrated uncertainty with missing-evidence flags,
inspection-as-noisy-observation, rare-event outcome treatment, a
mechanistic consequence surrogate, and decision-theoretic inspection
prioritisation in a single framework. That integration is the novelty claim,
and Section~\ref{sec:results} tests its components rather than asserting them.
